\scanpage{83}
$K^{\bullet\bullet}$ Algèbre de DR à p.d. élémentaire

a) Anneau bigradué \inserted{à bidegr. $\geq 0$}, \deleted{\unsure{comm.}} alterné pour \inserted{``deg. ext''}

b) \deleted{\unsure{Système}} Structure à p. d. sur l'idéal $\sum_{i+j>0}K^{ij}$

c) différentielle d'Anneau gradué (pour la grad ext) ne changeant pas le degré total, et comp. \unsure{avec|aux} p. d.

d) $d:K^{1,0}=\Phi\longrightarrow K^{11}\unsure{\equiv|=}\Psi$ est un épi.

e) Posant $k=K^{00}$, $K^{**}$ est l'algèbre \inserted{alternée} à p. d. libre engendrée par $\Phi$ et $\Psi$ i.e. $K^{**}\simeq\Gamma^*\Phi\otimes\Lambda^*\Psi$ [Il en résulte que $K^{ij}=0$ si $j>i$.]

f) \deleted{\unsure{$\Psi$ est $k$-plat|Mono-plat}} $\Phi$ et $\Psi$ sont $k$-plats (donc aussi $L=\operatorname{Ker}d=\mathscr H^{10}(K^{\bullet\bullet})$)

Il résulte déjà de ceci que
\[
H^{*i}(K^{\bullet\bullet})=
\begin{cases}
\Gamma^*L & \text{si }i=0\\
0 & \text{si }i>0
\end{cases}
\]
donc que $K^{\bullet\bullet}$ est une résolution de $\Gamma^*L$. On va supposer le plus souvent

g) $L$ est un $k$-Module \deleted{inversible}\inserted{loc. libre de rg} $1$.

(voire libre de rang $1$ et trivialisé ...)
% archive sheet number: 82
